% SIAM News Software and Programming — author draft, not submitted.
% Canonical editorial prose: article/feature.md. Keep both files synchronized.
% From article/: pdflatex -interaction=nonstopmode main.tex (twice)
% This is a magazine-style proof, NOT the official SIAM research-journal template.
\documentclass[11pt,letterpaper]{article}
\usepackage[T1]{fontenc}
\usepackage[utf8]{inputenc}
\usepackage[margin=23mm]{geometry}
\usepackage{lmodern}
\usepackage{microtype}
\usepackage{graphicx}
\usepackage{amsmath}
\usepackage{xcolor}
\usepackage{parskip}
\usepackage{hyperref}
\usepackage{caption}
\definecolor{ProxTeal}{HTML}{166B75}
\definecolor{ProxInk}{HTML}{253C49}
\definecolor{ProxGray}{HTML}{60717B}
\hypersetup{colorlinks=true,urlcolor=ProxTeal,linkcolor=ProxTeal,pdfauthor={Hector Maravillo; Diego Villarreal de la Cerda; Heriberto Espino-Montelongo},pdftitle={Who Counts as a Neighbor?}}
\captionsetup{font=small,labelfont={bf,color=ProxInk},justification=raggedright,singlelinecheck=false}
\renewcommand{\familydefault}{\sfdefault}
\renewcommand{\figurename}{Figure}
\setlength{\parindent}{0pt}
\setlength{\parskip}{0.48em}
\raggedbottom
\newcommand{\featurefigure}[3]{%
  \begin{figure}[htbp]
  \centering
  \IfFileExists{../site/assets/#1.png}{\includegraphics[width=.80\linewidth]{../site/assets/#1.png}}{%
  \fbox{\parbox[c][37mm][c]{.77\linewidth}{\centering\small Figure generated by the ProximityGraphs companion scripts\\[3pt]#1.png}}}%
  \caption{#2 \href{#3}{Explore the interactive figure}. Original static geometry and image courtesy of the authors; all sites are synthetic.}
  \end{figure}%
}
\begin{document}
\begin{center}
{\footnotesize\bfseries\color{ProxTeal} SOFTWARE AND PROGRAMMING \quad \textbullet \quad SIAM NEWS PROPOSAL}\\[8pt]
{\LARGE\bfseries\color{ProxInk} Who Counts as a Neighbor?}\\[5pt]
{\large\color{ProxGray} How ProximityGraphs turns one point cloud into many mathematical worlds}\\[9pt]
{\normalsize H\'ector Maravillo \quad Diego Villarreal de la Cerda \quad Heriberto Espino-Montelongo}\\[4pt]
{\scriptsize\color{ProxGray} Proposed byline; affiliations, author order, and biographies require contributor approval.}
\end{center}
\vspace{4pt}\hrule\vspace{9pt}

A planner studying a map of water stations wants to know which station serves each home. An ecologist wants to trace the outer boundary of scattered observations. A network engineer wants to decide which sensors might reasonably be connected. All three begin with points in the plane, but their questions require different meanings of \emph{neighbor}.

That distinction matters. Distance alone does not tell us which connections should exist: a pair of points may share a nearest-site boundary, belong to the same geometric envelope, or qualify as neighbors because a particular region between them contains no other point. Each choice produces a different mathematical structure---and potentially a different scientific interpretation.

\href{https://github.com/ProximityGraphs/ProximityGraphs}{ProximityGraphs} is an open-source Python library that makes these choices explicit. It brings together classic constructions such as Voronoi and Delaunay, Gabriel and relative neighborhood graphs, beta-skeletons, alpha-hulls, nearest-neighbor graphs, and related families under a common point-set workflow. Instead of coding every geometric rule from scratch, users can build, visualize, and compare alternatives on the same data.

\section*{A city divided by distance}

Consider a hypothetical city with temporary drinking-water stations. If people always visit their closest station \emph{as the crow flies}, each station occupies a Voronoi cell: the set of locations closer to it than to any other station. Formally, for sites \(P=\{p_1,\ldots,p_n\}\), the cell of site \(p_i\) is
\[
V_i=\{x:\|x-p_i\|\leq\|x-p_j\| \text{ for every } j\}.
\]
Connecting sites whose Voronoi cells share a boundary gives the Delaunay dual. The first object answers a coverage question; the second reveals which service territories meet. Both emerge from the same coordinates.

There is a compelling historical connection. During London's 1854 cholera outbreak, physician John Snow investigated the geographic relationship between deaths and water pumps. Later spatial analyses revisited nearest-pump catchments. Our diagram is \textbf{not} Snow's original map, and straight-line Voronoi regions are not a substitute for realistic travel-time, capacity, or demand models. They are a transparent starting point for explaining spatial allocation.

\featurefigure{figure01-water}{Hypothetical water stations with Euclidean Voronoi service cells and Delaunay-adjacency edges. Select or move a station to see which regions change. The library constructs the site-adjacency graph; finite Voronoi cells are clipped separately.}{https://heri-espino.github.io/proximitygraphs-siam-figures/\#water}

\section*{A boundary made from circles}

Now imagine ecological observations distributed around an irregular landscape. The familiar convex hull---like a rubber band stretched around the outermost points---ignores bays and indentations. It may enclose regions with no observations and conceal the outline's geometry.

The alpha-shape family introduces a scale into boundary construction. Its circular-arc relative, the \emph{alpha-hull}, can trace a boundary with arcs rather than straight segments. In the ProximityGraphs convention used in our example, an arc's supporting circle has radius \(R=1/|\alpha|\), with a negative alpha parameter. Changing this radius changes which concavities become visible.

The second visualization makes those circles tangible: select an arc to reveal its supporting circle, or move an observation and watch the envelope change. The outcome is a \textbf{geometric description of sampled points}, not a statistical estimate of habitat occupancy. Whether an outline represents a real ecological boundary depends on survey design and additional modeling.

\featurefigure{figure02-habitat}{Alpha-hull arcs around synthetic habitat observations, highlighting one supporting circle. The interactive version lets readers adjust \(\alpha\) and the point positions; original static arcs come from ProximityGraphs.}{https://heri-espino.github.io/proximitygraphs-siam-figures/\#habitat}

\section*{When an empty region creates an edge}

A third problem arises when points represent wireless sensors. Connecting every pair yields a dense network that says little about local structure. Connecting only the nearest pairs imposes a different rule. A \emph{beta-skeleton} provides an interpretable alternative: connect two sites only when a region determined by that pair is empty of other sites.

In the lune-based family, the relevant region is a lens formed by two overlapping disks. The parameter \(\beta\) controls its shape. At \(\beta=1\), the construction gives the Gabriel graph; at \(\beta=2\), the relative neighborhood graph, under standard general-position conventions. A point added inside the relevant lens can eliminate a previously permissible edge.

That local cause-and-effect is the focus of the third visualization. Readers can change \(\beta\), inspect a candidate edge, and see the two disks whose intersection must be empty. The graph identifies \emph{geometric candidates}, not guaranteed radio connections: physical deployment also requires propagation, terrain, power, and reliability models.

\featurefigure{figure03-sensors}{A beta-skeleton over hypothetical sensor locations with the selected edge's empty lens shown. Moving or adding a point can change the edge set. The static graph was generated with ProximityGraphs.}{https://heri-espino.github.io/proximitygraphs-siam-figures/\#sensors}

\section*{One point set, several answers}

These examples suggest a practical question for researchers: \textbf{What changes when we change the neighborhood rule but keep the data fixed?} ProximityGraphs is designed for precisely that comparison. Its point-set objects, graph constructors, visualization tools, and experiment helpers create a workflow in which a hypothesis about spatial structure can be tested rather than merely illustrated.

A small Python example captures the idea:
\begin{quote}
\ttfamily\small
import proximitygraphs as pg\\
points = pg.SetPoints.grid(shape=(3, 3))\\
gabriel = pg.GG(points)\\
relative = pg.RNG(points)\\
delaunay = pg.DelaunayG(points)
\end{quote}

The same nine sites produce three graph constructions that encode different definitions of adjacency. The library also provides utilities for repeated experiments and optional geospatial-data integration. It complements general graph and scientific-computing packages by focusing on the geometry that determines which edges exist.

To let readers explore the distinctions without installing Python, we built an \href{https://heri-espino.github.io/proximitygraphs-siam-figures/}{interactive companion}. Its original examples use geometry exported by ProximityGraphs; when readers add, delete, or drag points, independent browser code recomputes the demonstrations and labels them accordingly. The source, figures, and generation scripts are public and reproducible.

\textbf{A graph begins with a definition, not a drawing.} Before treating any line as a road, relationship, or possible communication link, ask which rule created it---and which scientific assumptions make that rule useful. ProximityGraphs provides a place to ask the question with code, and the interactive diagrams let anyone see the answer change.

\medskip
\textbf{Explore:} \href{https://github.com/ProximityGraphs/ProximityGraphs}{Source and Python package} \textbullet{}
\href{https://pypi.org/project/proximitygraphs/}{Installation} \textbullet{}
\href{https://heri-espino.github.io/proximitygraphs-siam-figures/}{Interactive figures}.
\quad \textbf{Install:} \texttt{pip install proximitygraphs}. (MIT License.)

\section*{References}
\begin{thebibliography}{6}
\bibitem{edelsbrunner} Edelsbrunner, H., Kirkpatrick, D. G., \& Seidel, R. (1983). On the shape of a set of points in the plane. \emph{IEEE Transactions on Information Theory, 29}(4), 551--559.
\bibitem{gabriel} Gabriel, K. R., \& Sokal, R. R. (1969). A new statistical approach to geographic variation analysis. \emph{Systematic Zoology, 18}(3), 259--278.
\bibitem{proximitygraphs} Maravillo, H., Villarreal de la Cerda, D., \& Espino-Montelongo, H. (2026). \emph{ProximityGraphs} [Computer software]. Zenodo. \url{https://doi.org/10.5281/zenodo.21088183}.
\bibitem{snow} Snow, J. (1855). \emph{On the mode of communication of cholera} (2nd ed.). John Churchill.
\bibitem{toussaint} Toussaint, G. T. (1980). The relative neighbourhood graph of a finite planar set. \emph{Pattern Recognition, 12}(4), 261--268.
\end{thebibliography}

\section*{About the authors}
\textbf{H\'ector Maravillo} develops mathematical and computational research involving proximity graphs. \textit{Confirm affiliation and biography.}

\textbf{Diego Villarreal de la Cerda} contributes to ProximityGraphs and its computational development. \textit{Confirm affiliation and biography.}

\textbf{Heriberto Espino-Montelongo} studies actuarial science and data science at Universidad de las Am\'ericas Puebla and works on geometric and computational methods. \textit{Confirm preferred biography.}

{\footnotesize\color{ProxGray}Editorial draft. All demonstrations use synthetic points. The authors must approve this text and the figure credits before any SIAM News submission.}
\end{document}
